% ============================================================
% Apartado 2: Marco Empírico (Reglamento de Titulación, Art. 25 y 27)
% ============================================================
\chapter{Marco Empírico}
\label{cap:marco-empirico}

% Art. 27: el marco empírico debe describir (a) brevemente la empresa o
% institución donde se analiza o aplica el tema; (b) el análisis estratégico que
% identifica el problema y sus causas, junto con los resultados de investigación
% obtenidos de fuentes primarias; y (c) las soluciones propuestas por el alumno.

\section{Descripción de la institución}

Breve presentación de la empresa u organización donde se desarrolla el estudio
(reseña, rubro, estructura, situación actual).

\section{Metodología de la investigación}

\subsection{Tipo y diseño de investigación}
Enfoque (cuantitativo, cualitativo, mixto), tipo y diseño.

\subsection{Población y muestra}
Definición de la población y criterios de selección de la muestra.

\subsection{Técnicas e instrumentos de recolección de datos}
Encuestas, entrevistas, observación, análisis documental, etc.

\subsection{Metodología de desarrollo del sistema}
Si el trabajo incluye un software, metodología empleada (Scrum, RUP, XP,
cascada, etc.) y sus fases.

\section{Análisis estratégico}

Identificación del problema y sus causas a partir de la situación actual; uso de
herramientas de análisis (FODA, diagramas de causa-efecto, etc.).

\section{Resultados de la investigación}

Presentación de los resultados obtenidos de fuentes primarias, apoyada en tablas
y figuras (numeradas según APA).

% --- Ejemplo de figura ---
% \begin{figure}[H]
%   \centering
%   \includegraphics[width=0.8\textwidth]{nombre-imagen}
%   \caption{Descripción de la figura.}
%   \label{fig:ejemplo}
% \end{figure}

% --- Ejemplo de tabla (estilo APA) ---
\begin{table}[H]
  \centering
  \caption{Ejemplo de tabla de resultados.}
  \label{tab:ejemplo}
  \begin{tabular}{lcc}
    \toprule
    \textbf{Variable} & \textbf{Valor} & \textbf{Porcentaje} \\
    \midrule
    Categoría A & 120 & 60\,\% \\
    Categoría B & 80  & 40\,\% \\
    \midrule
    \textbf{Total} & \textbf{200} & \textbf{100\,\%} \\
    \bottomrule
  \end{tabular}
\end{table}

\section{Propuesta de solución}

Soluciones propuestas por el autor. Si el trabajo desarrolla un sistema:

\subsection{Análisis del sistema}
Requerimientos funcionales y no funcionales; casos de uso.

\subsection{Diseño del sistema}
Arquitectura, modelo de datos y diseño de interfaces.

\subsection{Implementación}
Tecnologías utilizadas y aspectos relevantes de la construcción.

\begin{lstlisting}[language=Python,caption={Ejemplo de fragmento de código}]
def saludar(nombre):
    return f"Hola, {nombre}"
\end{lstlisting}

\subsection{Pruebas}
Estrategia y resultados de las pruebas realizadas.
